\paragraph{A local inverse response on the interacting infinite ground.}
The original zero-chemical-potential Weyl/contact action on $\mathbb Z^3$
admits a controlled spatial inverse response on the earlier
constructed thermodynamic ground, at fixed spacing and away from real
frequency.  This does not extrapolate the special two-eigenvalue cube
identity.  Let $\omega$ be the translation-averaged full-ground limit,
$L_\omega$ its positive GNS generator, and
\[
L=L_\omega\oplus-\overline{L_\omega},\qquad
V_{x\sigma}=c_{x\sigma}^\dagger\Omega
                 \oplus\overline{c_{x\sigma}\Omega}.
\]
The CAR give an isometry $V$ from
$\ell^2(\mathbb Z^3;\mathbb C^2)$.  Define $J$ in the same way from
$b_{x\sigma}=n_{x\bar\sigma}c_{x\sigma}$ and write
$\rho_x{}_{\sigma\tau}=\omega(c_{x\tau}^\dagger c_{x\sigma})$.
Exact onsite CAR identities, with vanishing different-site
anticommutators, give
\begin{align*}
R&=\bigoplus_x\bigl(\operatorname{tr}(\rho_x)I-\rho_x\bigr),
 &V^\dagger J&=J^\dagger V=J^\dagger J=R,\\
LV&=Vh+UJ,
 &0&\leq R\leq I,\quad \|LV\|\leq\sqrt3c+U.
\end{align*}
Finite-support approximants and closedness of $L$ extend the last
identity to the entire one-particle input space.  With $P=VV^\dagger$,
$Q=I-P$, the bounded off-diagonal coupling is $W=QLV=U(J-VR)$.
Subtracting $WV^\dagger+VW^\dagger$ from $L$ proves self-adjointness
of its compressed complement $T=QLQ|_{Q\operatorname{Dom}L}$.
The exact infinite-space Schur identity is therefore
\begin{align*}
\mathcal G(z)&=V^\dagger(z-L)^{-1}V
          =[zI-h-UR-\Sigma(z)]^{-1},\\
\Sigma(z)&=W^\dagger(z-T)^{-1}W,
 &W^\dagger W&=U^2(R-R^2).
\end{align*}
In particular $\|\Sigma(z)\|\leq U^2/(4\operatorname{Im}z)$.

This ground has strictly nonzero memory for $U>0$.  The source hidden-spin
frame obeys $R_{x+t}=\pm R_tR_x$, so translations normalize its global
$\mathrm{SU}(2)$ symmetry.  Every full-box ground projector and the
exterior vacuum are invariant; translated averages and weak-star limits
retain the symmetry.  Hence $\rho_x=rI$, with the earlier density theorem
giving $1/8\leq r\leq1/2$, and
\[
 \frac{7U^2}{64}I\leq W^\dagger W
       =U^2r(1-r)I\leq\frac{U^2}{4}I.
\]
The spectral theorem gives $iy\Sigma(iy)\to W^\dagger W$ strongly as
$y\to\infty$.  Thus the same-mode response cannot be that of a
number-conserving static quadratic Hamiltonian with zero memory.
Neither a scattering theorem
nor faithfulness of the local ground-state density matrices follows.

\paragraph{A quantitative connected four-point witness.}
The same ground has a non-Gaussian multi-time field, a stronger statement
than nonzero two-point memory.  Fix $a=c_{x\sigma}$ and
$n=n_{x\bar\sigma}$, and set
\[
 A(t)=\{\tau_t(a),a^\dagger\},\qquad
 D(t)=A(t)-\omega(A(t))I,\qquad
 F(t)=\omega(D(t)^\dagger D(t)).
\]
The local commutator gives $D(0)=0$ and
$D'(0)=-iU(n-rI)$.  Since $n^2=n$,
$F(t)/t^2\to U^2r(1-r)\geq7U^2/64$.
An explicit remainder uses
$M=9c^2+12cU+U^2$ and $\delta=[H,\cdot]$:
\[
 \|\delta^2(a)\|
 \leq3c(3c+U)+U(6c+3c+U)=M,
 \qquad \|D(t)+itU(n-rI)\|\leq2Mt^2.
\]
Thus, for $U>0$,
\[
 F(t)\geq\frac{U^2t^2}{64}>0,
 \qquad 0<|t|\leq\frac{U}{16M}.
\]
For odd zero-mean operators define the connected fermionic function
\begin{align*}
 \kappa_4(B_1,B_2,B_3,B_4)
 &=\omega(B_1B_2B_3B_4)
   -\omega(B_1B_2)\omega(B_3B_4)\\
 &\quad+\omega(B_1B_3)\omega(B_2B_4)
   -\omega(B_1B_4)\omega(B_2B_3).
\end{align*}
With $(B_1,B_2,B_3,B_4)=(a,\tau_t(a)^\dagger,\tau_t(a),a^\dagger)$,
$F(t)$ is the sum of $\kappa_4$ on words $1234,1243,2134,2143$.
All crossed contractions, including anomalous ones, cancel pairwise.
At least one connected function therefore has magnitude
$U^2t^2/256$ or greater on the displayed interval.  This excludes a
quasifree representation of the complete multi-time field by linear
fermionic fields, including pairing and auxiliary Gaussian modes.
Indeed any quadratic evolution on an enlarged CAR algebra preserves
linearity of an observed mode, making its anticommutator scalar and
$D(t)=0$ in every reference state.  Thus even a non-Gaussian reference
cannot realize this witness through such a linear quadratic embedding.
It does not exclude nonlinear encodings or establish continuum scattering.

\paragraph{The full ordered Euclidean hierarchy.}
The same positive ground generator $L_\omega\geq0$ defines, for ordered
real times $s_1\leq\cdots\leq s_n$ and bounded local CAR operators,
\[
 S(A_1,s_1;\ldots;A_n,s_n)
 =\langle\Omega,A_1e^{-(s_2-s_1)L_\omega}A_2\cdots
                   e^{-(s_n-s_{n-1})L_\omega}A_n\Omega\rangle,
 \qquad |S|\leq\prod_j\|A_j\|.
\]
Every finite approximant uses one common translated full-ground state
and its positive generator $H_\Lambda-E_\Lambda$ for the whole word,
with zero exterior Hamiltonian, before averaging over translations.
For each positive gap the spectral theorem gives
\[
 e^{-sL}=\int_{\mathbb R}\frac{s}{\pi(s^2+t^2)}e^{itL}\,dt,
 \qquad L\geq0.
\]
The product integral is dominated by $\prod_j\|A_j\|$.
Its real-time integrand is the expectation of
$A_1\tau_{t_1}(A_2)\tau_{t_1+t_2}(A_3)\cdots$.
Local-dynamics bounds, telescoping the product, and the same
weak-star ground subsequence give uniform convergence on compact time
boxes.  The omitted product-measure tail is at most
$\sum_j2s_j/(\pi T_j)$, capped by one; a finite-versus-limit comparison
costs twice this tail times the product of operator norms.  Dominated
convergence therefore supplies every ordered Euclidean $n$-point limit
on this one ground, without another subsequence or a state-convergence rate.
Zero gaps are handled by the identity operator.

For a positive-time word $F=(A_1,t_1;\ldots;A_n,t_n)$ define
\[
 \Psi_F=e^{-t_1L_\omega}A_1e^{-(t_2-t_1)L_\omega}A_2\cdots
                      e^{-(t_n-t_{n-1})L_\omega}A_n\Omega.
\]
The reflected list is
$(A_n^*,-t_n;\ldots;A_1^*,-t_1)$, in that order.
Concatenation with a second positive-time word gives
$S(\theta F,G)=\langle\Psi_F,\Psi_G\rangle$; hence every reflected
correlation matrix is positive semidefinite.  No CAR factors are
reordered.  Exact rational two-mode controls check the Gram identity,
all 31 principal minors of a five-word example, and a negative odd-word
cross entry.  This proves ordered time-reflection positivity at fixed
spacing.  It does not establish $O(4)$ covariance, bosonic permutation
symmetry, a spatial continuum, or the full relativistic Euclidean axioms.

\paragraph{Uniform inverse and spatial decay bounds.}
Fix $\operatorname{Im}z\geq\eta>0$, $|z|\leq Z$, with $Z\geq\eta$.
For a unit vector $v=Vf$, Cauchy--Schwarz applied to
$\langle(\bar z-L)v,(z-L)^{-1}v\rangle=1$ yields
\[
 -\operatorname{Im}\mathcal G(z)\geq B_0^{-1}I,
 \qquad
 \|\mathcal G(z)^{-1}\|\leq B_0,
 \qquad B_0=\frac{(Z+2c+U)^2}{\eta}.
\]
Both $\mathcal G$ and its adjoint are bounded below; its range is
therefore closed and dense, which proves surjectivity in infinite volume.
No determinant or finite-volume gap is used.

Let $d=\|x-y\|_\infty$.  The original graded locality bound gives
\begin{align*}
 \|\mathcal G_{xy}(z)\|&\leq A(1+d)^2e^{-2\mu d},\\
 \mu&=\min\{1/8,\eta/(384c)\},
 &A&=192c/\eta^2+8/\eta.
\end{align*}
For $d\geq2$, approximate the evolved odd operator in the radius
$d-1$ cube, whose anticommutator with the remote mode is zero.
Split the Laplace integral at $(d-1)/(192c)$ and use the previous
full-CAR boundary estimate on the first part and the uniform norm bound
on its tail.  This proves the displayed constants without assuming
spatial clustering of $\omega$.

Put $q=(1+\mu)^{-1}$ and
$T_n(q)=\sum_{d\geq1}d^nq^d$.  These are rational functions with
positive Eulerian coefficients.  Define
\begin{align*}
 \widehat S&=24T_5(q)+48T_4(q)+26T_3(q)+4T_2(q)+2T_1(q),\\
 \alpha&=\min\{\mu,(2B_0A\widehat S)^{-1}\}>0.
\end{align*}
For any bounded real one-Lipschitz function $f$, the block Schur test gives
\[
 \|e^{\alpha f}\mathcal G e^{-\alpha f}-\mathcal G\|
       \leq\alpha A\widehat S\leq(2B_0)^{-1}.
\]
A Neumann series bounds the conjugated inverse by $2B_0$.
Choosing $f(u)=\min\{\|u-y\|_\infty,R\}$ with $R\geq d$ proves
\begin{align*}
 \|(\mathcal G^{-1})_{xy}(z)\|&\leq2B_0e^{-\alpha d},\\
 \|\Sigma_{xy}(z)\|&\leq(Z+U+2B_0+2c)e^{-\alpha d}.
\end{align*}
Only bounded weights were used.  The constants are conservative bounds,
not estimates of a physical correlation length.

\paragraph{The controlled spatial symbol and its geometric boundary.}
Translation invariance and exponential summability give an analytic
two-by-two inverse-response symbol $K(k,z)$ for
$\|\operatorname{Im}k\|_1<\alpha$.
For real momentum increment $p$, its total-degree-$n$ Taylor remainder
is bounded by
\[
 \frac{2B_0\|p\|_1^{n+1}}{(n+1)!}
 \sum_{x\in\mathbb Z^3}\|x\|_\infty^{n+1}
                       e^{-\alpha\|x\|_\infty}.
\]
The shell count $24d^2+2$ and the same rational $T_n$ functions give
explicit upper bounds on this convergent sum.  A frequency circle of
radius $\eta/2$ gives, for $|w|<\eta/2$, the additional remainder
$B_* (2|w|/\eta)^{n+1}/(1-2|w|/\eta)$, where $B_*$ is the above
inverse bound evaluated at $(\eta/2,Z+\eta/2)$.

These estimates supply actual local spatial coefficients on the
interacting infinite ground, with a controlled complex-frequency
domain.  They do not compute the ground functional or those coefficients,
and deteriorate as $\eta\to0$.  A real hyperbolic principal symbol,
common matter coefficients and full action normalization are
required before applying the existing \texttt{wave\_density\_metric}
and \texttt{source\_tetrad} interfaces.  The joint modular, area,
stress and Einstein balance requirements remain separate.

For a refinement $c=v/a$ at fixed $v,U,Z,\eta>0$, the retained estimates
have $\alpha=O(a^9)$ and therefore physical momentum width
$\alpha/a=O(a^8)$; the connected-witness time window is $O(a^2)$.
These shrinking bounds do not prove uniform continuum control or a
no-go theorem for the action.
